\chapter{Discussion}
\label{chap:discussion}

This chapter interprets the results of Chapter~\ref{chap:results} against the literature of Chapter~\ref{chap:lit} and the research questions. Its main limitation is the one stated throughout: every result is from the bench.

\section{Deterministic acquisition and what it costs}

The synthesis of the literature review noted that edge-telemetry studies are often judged by algorithm accuracy or mean latency, without jitter, queue behaviour or outage recovery (Chapter~\ref{chap:lit}). The results here report those quantities directly: zero dropped samples, the worst gap with and without logging, and the cause of the only recurring gap. The flash-write pause is the clearest example of why they matter. Its rate was cut sixfold by buffering, but its length was not, because the pause comes from the flash itself. A 40\,ms gap in a 200\,Hz stream is eight missing sample instants once a minute, which is small for 1\,s ride windows but would matter for impact detection. Moving logging to separate storage, or to the Raspberry Pi, would be expected to remove it rather than reduce it.

\section{Ride characterisation}

The calibration result supports the review's warning that low-cost inertial measurements depend on the individual sensor: two nominally identical units disagreed by about 8\% at rest, within their datasheet tolerance, and their correction factors changed when they were moved between boards. A ride index that is to be compared across sessions therefore has to be calibrated in its final mounting, which the firmware supports with a single command. The features themselves are computed correctly, but they use the acceleration magnitude rather than a vehicle-fixed vertical axis and have no ISO~2631 frequency weighting, so they describe overall vibration exposure rather than ride comfort as the standard defines it. These are deliberate first-pass simplifications, and both can be added in software once vehicle data exist. The sensors' $\pm8g$ range and 260\,Hz filter are also provisional: the filter lets content between 100 and 260\,Hz alias into the measured band, and the range should be set from the peaks the vehicle actually produces.

RQ4 cannot be answered from the bench. The method is in place: the motor's Hall pulses give its speed, so vibration that tracks motor speed with the vehicle stationary on stands can be separated from vibration that appears only on the road. That separation depends on the Hall tap being wired and on the first vehicle session.

\section{Hazard sensing}

The camera pipeline is fast enough for its purpose: an alert about 0.2\,s after an object appears is about 1.7\,m of travel at 30\,km/h, small against the stopping distance at that speed. The weak point is distance, not speed. The first fused hazard put a person at 0.874\,m by the camera's width heuristic and 0.096\,m by the time-of-flight sensor, so a camera-only distance band could be wrong by a whole band at close range. This is the case for fusion made in Section~\ref{sec:camera-tof-fusion}: the camera says what and in which direction, the sensor says how far. The time-of-flight range of 1.2\,m, however, is only about 0.14\,s of travel at 30\,km/h, so the stand-in sensors serve low-speed manoeuvring and the camera remains the only source at road speed until longer-range sensing is fitted.

The pothole models show the same dependence on local data that the review noted for road-roughness classifiers, whose reported accuracies do not transfer without retraining (Section~\ref{sec:lit-signal-processing}). A model that scored 59.7\% on test images from its own source found almost none of the potholes in South African dashcam footage, and training on tiles cut from the road band of that footage raised the validation score to 0.714. The held-out test score is needed before this improvement can be claimed, and the model has not yet been run on the vehicle's camera, whose view, resolution and colour, without an infrared filter, differ from the dataset's.

\section{Data integrity and connectivity}

Logging independently of the connection, and reconnecting when the modem reports that data are available, gave the store-and-forward behaviour the review identified as necessary for intermittent coverage. The read-back is the important result: it recovered every record, and it found a fault, invalid JSON for missing values, that no amount of watching the live stream had revealed. Two parts of the cellular path remain open: the data go to a public test broker rather than a database, and the backlog of unsent records is not yet replayed after an outage.

\section{Driver display}

The decision to present on the vehicle's existing dashboard rather than a windshield HUD is supported by the review: a head-up position does not by itself make a display safe, its optics are hard to make legible on a removable unit in daylight, and a second display adds a competing channel (Section~\ref{sec:dashboard-page}). Two results bear directly on the safety of what is shown. Live-only mode removes a real hazard found during integration: the dashboard had been showing simulated speed and battery values without any sign when the data link was down. And the touch measurement shows that the panel is not a source of delay, so the lag noticed in bench use comes from the buttons acting on release, which can be changed. What RQ2 asks, whether the display can be read within the 2\,s guidance, has not been measured; the figures for character size, contrast and touch-target size from Chapter~\ref{chap:lit} give a check that can be applied before the on-vehicle observation.

\section{Integration and power}

The most consequential finding of the integration work was not in any one subsystem. The telemetry unit's LTE modem, powered from the Raspberry Pi's USB port, pulled the supply into under-voltage whenever it transmitted, throttled the processor and slowed the detector, and is the likely reason the Raspberry Pi~4 stopped booting on 6 October. The modem worked correctly throughout; the fault was in how the parts were powered together. A separate supply for the telemetry unit is therefore a design requirement for the vehicle installation, not an option.

\section{Validity and limitations}

All results come from one set of hardware on the bench, with the vehicle's vibration, heat, supply and lighting absent. Several parameters are still defaults or uncalibrated: the camera's focal length, the time-of-flight sensor directions, the accelerometer range and filter, and the GNSS speed units. The pothole score is on validation tiles. These limits do not invalidate the bench results, which measure what they claim to measure, but they bound what can be concluded from them, as the next chapter sets out.
